%%%%%%%%%%%%%%%%%%%%%%%%%%%%%%%%%%%%%%%%%%%%%%%%%%%%%%%%%%%%%%%%%%%%%%%%%%%%%%%
%% Pacotes do documento
%%
%% Os do modelo iCEV (indentfirst, abntex2cite) e os que o conteúdo
%% exige (siunitx, glossaries, nomencl, booktabs). A ordem importa:
%%   glossaries depois do hyperref, que a abntex2 carrega.
%%%%%%%%%%%%%%%%%%%%%%%%%%%%%%%%%%%%%%%%%%%%%%%%%%%%%%%%%%%%%%%%%%%%%%%%%%%%%%%

%% --- Fontes e codificação --------------------------------------------------
\usepackage[T1]{fontenc}
\usepackage[utf8]{inputenc}
\usepackage{indentfirst}
%% Times no corpo, como no modelo. O modelo usa mathptmx; newtx compõe a mesma
%% Times com matemática completa (negrito, letras gregas), que os capítulos
%% usam. As figuras do pipeline (experiments/figures/scripts/theme_kaelix.R)
%% compõem em Times New Roman para casar com este corpo.
\usepackage{newtxtext,newtxmath}
\usepackage{microtype}

%% --- Pacotes de conteúdo ---------------------------------------------------
\usepackage{graphicx}
\usepackage{xcolor}      % marcas vermelhas \emandamento; substitui o color do modelo
\usepackage{booktabs}    % réguas de tabela
\usepackage{amsmath}
\usepackage{caption}
%% Título antes da ilustração ou da tabela (NBR 14724, 5.8 e 5.9).
\captionsetup{justification=justified,singlelinecheck=false}
\captionsetup[figure]{position=top}
\captionsetup[table]{position=top}
\usepackage{letltxmacro} % \LetLtxMacro para o \fonte de estilo-icev.tex
\usepackage{siunitx}
%% Separador decimal do português: vírgula. Sem isto, \num{0,5} sai "0.5".
%% Incerteza: \SI{75,0 \pm 3,5}{\percent} sai "(75,0 ± 3,5) %", com a unidade
%% valendo para os dois números; para "75,0 % ± 3,5 %", troque bracket por repeat.
\sisetup{
  output-decimal-marker      = {,},
  group-separator            = {\,},
  range-phrase               = {~a~},
  list-final-separator       = {~e~},
  per-mode                   = symbol,
  uncertainty-mode           = separate,
  separate-uncertainty-units = bracket
}

%% --- Citações (ABNT, sistema autor-data) -----------------------------------
\usepackage[alf,bibleftalign]{abntex2cite}
%% Estilo local, ajustado à NBR 6023:2025 e à NBR 10520:2023; as alterações
%% estão listadas no topo do arquivo. O abntex2-alf da distribuição segue as
%% normas de 2018 e 2002 e não imprime DOI.
\bibliographystyle{abntex2-alf-2025}
%% A NBR 6023 não usa < > em volta do endereço eletrônico desde 2018; o
%% abntex2cite ainda os acrescenta a todo \url.
\def\UrlLeft{}
\def\UrlRight{}

%% --- Listas de siglas e de símbolos ----------------------------------------
%% Sem as opções que põem as listas no sumário: elementos pré-textuais não podem
%% constar nele (NBR 6027, 6.3).
\usepackage[acronym,nomain,nonumberlist]{glossaries}
\makeglossaries
\usepackage{nomencl}
\makenomenclature

%% --- Caminhos para o repositório -------------------------------------------
%% Definidos uma única vez aqui; os capítulos têm fallback idêntico.
\newcommand{\kxfig}[1]{../../experiments/figures/output/#1}
\newcommand{\kxhw}[1]{../../#1}
